
\section{Hamiltonian Monte Carlo}

Hamiltonian dynamics is an alternative way to formulate Newtonian
mechanics. The Hamiltonian $H$ captures both the potential and kinetic
energy of a particle as a function of its position and velocity. The
dynamics can be described by the following differential equations.
\begin{align*}
\frac{dx}{dt} & =\frac{\partial H(x,v)}{\partial v},\\
\frac{dv}{dt} & =-\frac{\partial H(x,v)}{\partial x}.
\end{align*}
These equations preserve the Hamiltonian function $H$. In the simplest
Euclidean setting, it can be defined as follows.
\[
H(x,v)=f(x)+\frac{1}{2}\norm v^{2}
\]
so that 
\[
\frac{dx}{dt}=v,\quad\frac{dv}{dt}=-\nabla f(x)
\]
or
\[
\frac{d^{2}x}{dt^{2}}=-\nabla f(x).
\]

More generally, the Hamiltonian can depend on a function that defines
a local metric:
\[
H(x,v)=f(x)+\frac{1}{2}\log((2\pi)^{n}\det g(x))+\frac{1}{2}v^{T}g(x)^{-1}v
\]
where $g(x)$ is a matrix, and when it is positive definite, it defines
a local norm at $x$. In this sense, we can view the dynamics as evolving
on a manifold with local metric $g(x)$. Here in this chapter, we
will focus on the case when $g(x)=I$, the standard Euclidean metric.

(Riemannian) Hamiltonian Monte Carlo (or RHMC) \cite{Neal1996,Neal2012}\cite{Girolami2009,Girolami2011}
is a Markov chain Monte Carlo method for sampling from a desired distribution.
Each step of the method consists of the following: At a current point
$x$,
\begin{enumerate}
\item Pick a random velocity $y$ according to a local distribution defined
by $x$ (in the simplest setting, this is the standard Gaussian distribution
for every $x$).
\item Move along the Hamiltonian curve defined by Hamiltonian dynamics at
$(x,y)$ for time (distance) $\delta$.
\end{enumerate}
For the choice of $H$ above, the marginal distribution of the current
point $x$ approaches the target distribution with density proportional
to $e^{-f}$. For the exact Hamiltonian flow $T_{\step}$, no Metropolis
correction is needed.In numerical HMC, an inexact reversible volume-preserving
integrator is usually paired with a Metropolis filter to correct discretization
error. For exact flow, there is no discretization step size to tune.
For numerical integrators, the integrator step size is still constrained
by accuracy and acceptance probability, even though the Metropolis
filter corrects the stationary distribution. Hamiltonian Monte Carlo
can be used for sampling from a general distribution $e^{-H(x,y)}$. 
\begin{defn}
Given a continuous, twice-differentiable function $H:\mathcal{M}\times\Rn\subset\R^{n}\times\R^{n}\rightarrow\R$
(\emph{Hamiltonian}) where $\mathcal{M}$ is the domain of the $x$-coordinate
of $H$, we say $(x(t),y(t))$ follows a \emph{Hamiltonian curve}
if it satisfies the \emph{Hamiltonian equations}

\begin{align}
\frac{dx}{dt} & =\frac{\partial H(x,y)}{\partial y},\nonumber \\
\frac{dy}{dt} & =-\frac{\partial H(x,y)}{\partial x}.\label{eq:ham}
\end{align}
We define the map $T_{\step}(x,y)\defeq(x(\step),y(\step))$ where
$(x(t),y(t))$ follows the Hamiltonian curve with the initial condition
$(x(0),y(0))=(x,y).$
\end{defn}

Hamiltonian Monte Carlo is a Markov chain generated by a sequence
of random Hamiltonian curves.

\begin{algorithm2e}

\caption{Hamiltonian Monte Carlo}

\SetAlgoLined

\textbf{Input:} some initial point $x^{(1)}\in\mathcal{M}$.

\For{$k=1,2,\cdots,T$}{

Sample $y$ according to $e^{-H(x^{(k)},y)}/\pi(x^{(k)})$ where $\pi(x)=\int_{\Rn}e^{-H(x,y)}dy$.

With probability $\frac{1}{2}$, set $(x^{(k+1)},y^{(k+1)})=T_{\step}(x^{(k)},y)$. 

Otherwise, $(x^{(k+1)},y^{(k+1)})=T_{-\step}(x^{(k)},y)$.

}

\textbf{Output:} $(x^{(T+1)},y^{(T+1)})$.

\end{algorithm2e}

\subsection*{Time-reversibility }
\begin{lem}[Energy Conservation]
\label{lem:ham_eng_pre}For any Hamiltonian curve $(x(t),y(t))$,
we have that
\[
\frac{d}{dt}H(x(t),y(t))=0.
\]
\end{lem}

\begin{proof}
Note that
\[
\frac{d}{dt}H(x(t),y(t))=\frac{\partial H}{\partial x}\frac{dx}{dt}+\frac{\partial H}{\partial y}\frac{dy}{dt}=\frac{\partial H}{\partial x}\frac{\partial H}{\partial y}-\frac{\partial H}{\partial y}\frac{\partial H}{\partial x}=0.
\]
\end{proof}
\begin{lem}[Measure Preservation]
\label{lem:ham_measure_pre}For any $t\geq0$, we have that
\[
\det\left(DT_{t}(x,y)\right)=1
\]
where $DT_{t}(x,y)$ is the Jacobian of the map $T_{t}$ at the point
$(x,y)$.
\end{lem}

\begin{proof}
Let $(x(t,s),y(t,s))$ be a family of Hamiltonian curves given by
$T_{t}(x+sd_{x},y+sd_{y})$. We write 
\[
u(t)=\frac{\partial}{\partial s}x(t,s)|_{s=0}\,,\,v(t)=\frac{\partial}{\partial s}y(t,s)|_{s=0}.
\]
By differentiating the Hamiltonian equations (\ref{eq:ham}) with
respect to $s$, we have that
\begin{align*}
\frac{du}{dt} & =\frac{\partial^{2}H(x,y)}{\partial y\partial x}u+\frac{\partial^{2}H(x,y)}{\partial y\partial y}v,\\
\frac{dv}{dt} & =-\frac{\partial^{2}H(x,y)}{\partial x\partial x}u-\frac{\partial^{2}H(x,y)}{\partial x\partial y}v,\\
(u(0),v(0)) & =(d_{x},d_{y}).
\end{align*}
This can be captured by the following matrix ODE
\begin{align*}
\frac{d\Phi}{dt} & =\left(\begin{array}{cc}
\frac{\partial^{2}H(x(t),y(t))}{\partial y\partial x} & \frac{\partial^{2}H(x(t),y(t))}{\partial y\partial y}\\
-\frac{\partial^{2}H(x(t),y(t))}{\partial x\partial x} & -\frac{\partial^{2}H(x(t),y(t))}{\partial x\partial y}
\end{array}\right)\Phi(t)\\
\Phi(0) & =I
\end{align*}
using the equation
\[
DT_{t}(x,y)\left(\begin{array}{c}
d_{x}\\
d_{y}
\end{array}\right)=\left(\begin{array}{c}
u(t)\\
v(t)
\end{array}\right)=\Phi(t)\left(\begin{array}{c}
d_{x}\\
d_{y}
\end{array}\right).
\]
Therefore, $DT_{t}(x,y)=\Phi(t)$. Next, we observe that
\[
\frac{d}{dt}\log\det\Phi(t)=\tr\left(\Phi(t)^{-1}\frac{d}{dt}\Phi(t)\right)=\tr\left(\begin{array}{cc}
\frac{\partial^{2}H(x(t),y(t))}{\partial y\partial x} & \frac{\partial^{2}H(x(t),y(t))}{\partial y\partial y}\\
-\frac{\partial^{2}H(x(t),y(t))}{\partial x\partial x} & -\frac{\partial^{2}H(x(t),y(t))}{\partial x\partial y}
\end{array}\right)=0.
\]
Hence,
\[
\det\Phi(t)=\det\Phi(0)=1.
\]
\end{proof}
Using the previous two lemmas, we can see that Hamiltonian Monte Carlo
indeed converges to the desired distribution. 
\begin{lem}[Time reversibility]
\label{lem:time_reversibility}Let $p_{x}(x')$ denote the probability
density of one step of the Hamiltonian Monte Carlo starting at $x$.
We have that\textup{
\[
\pi(x)p_{x}(x')=\pi(x')p_{x'}(x)
\]
for almost every $x$ and $x'$ where $\pi(x)=\int_{\Rn}e^{-H(x,y)}dy$.}
\end{lem}

\begin{proof}
Fix $x$ and $x'$. Let $F^{x}_{\step}(y)$ be the $x$ component
of $T_{\step}(x,y)$. Let $V_{+}=\{y:\text{ }F^{x}_{\step}(y)=x'\}$
and $V_{-}=\{y:\text{ }F^{x}_{-\step}(y)=x'\}$. Then, 
\[
\pi(x)p_{x}(x')=\frac{1}{2}\int_{y\in V_{+}}\frac{e^{-H(x,y)}}{\left|\det\left(DF^{x}_{\step}(y)\right)\right|}+\frac{1}{2}\int_{y\in V_{-}}\frac{e^{-H(x,y)}}{\left|\det\left(DF^{x}_{-\step}(y)\right)\right|}.
\]
We note that this formula assumes that $DF^{x}_{\step}$ is invertible.
Sard's theorem shows that $F^{x}_{\step}(N)$ has measure $0$ where
$N\defeq\{y:\,DF^{x}_{\step}(y)\text{ is not invertible}\}$. The
same argument applies to $F^{x}_{-\step}$. Therefore, the formula
is correct except for a measure zero subset of target points.

By reversing time for the Hamiltonian curve, we have that for the
same $V_{\pm}$, 
\begin{equation}
\pi(x')p_{x'}(x)=\frac{1}{2}\int_{y\in V_{+}}\frac{e^{-H(x',y')}}{\left|\det\left(DF^{x'}_{-\step}(y')\right)\right|}+\frac{1}{2}\int_{y\in V_{-}}\frac{e^{-H(x',y')}}{\left|\det\left(DF^{x'}_{\step}(y')\right)\right|}\label{eq:ham_transit}
\end{equation}
where $y'$ denotes the $y$ component of $T_{\step}(x,y)$ and $T_{-\step}(x,y)$
in the first and second terms, respectively.

We compare the first terms in both equations. Let $DT_{\step}(x,y)=\left(\begin{array}{cc}
A & B\\
C & D
\end{array}\right)$. At the regular points used in the change-of-variables formula, the
block$B=DF^{x}_{\step}(y)$ is invertible. Let $DT_{-\step}(x',y')=\left(\begin{array}{cc}
* & M\\*
 & *
\end{array}\right)$. Then $M=DF^{x'}_{-\step}(y')$. Since $T_{-\step}$ is the inverse
of $T_{\step}$, this matrix is the inverse of $DT_{\step}(x,y)$.
To compute $M$, solve $\left(\begin{array}{cc}
A & B\\
C & D
\end{array}\right)\binom{u}{v}=\binom{0}{q}$. The first block equation gives $v=-B^{-1}Au$, and the second gives
$(C-DB^{-1}A)u=q$. Thus $M=(C-DB^{-1}A)^{-1}$. Therefore, we have
that $DF^{x}_{\step}(y)=B$ and $DF^{x'}_{-\step}(y')=(C-DB^{-1}A)^{-1}$.
Hence, we have that
\begin{align*}
\left|\det\left(DF^{x'}_{-\step}(y')\right)\right| & =\left|\det(C-DB^{-1}A)^{-1}\right|=\frac{\left|\det B\right|}{\left|\det\left(\begin{array}{cc}
A & B\\
C & D
\end{array}\right)\right|}.
\end{align*}
Using that $\det\left(DT_{t}(x,y)\right)=\det\left(\begin{array}{cc}
A & B\\
C & D
\end{array}\right)=1$ (Lemma \ref{lem:ham_measure_pre}), we have that
\[
\left|\det\left(DF^{x'}_{-\step}(y')\right)\right|=\left|\det\left(DF^{x}_{\step}(y)\right)\right|.
\]
Hence, we have that
\begin{align*}
\frac{1}{2}\int_{y\in V_{+}}\frac{e^{-H(x,y)}}{\left|\det\left(DF^{x}_{\step}(y)\right)\right|} & =\frac{1}{2}\int_{y\in V_{+}}\frac{e^{-H(x,y)}}{\left|\det\left(DF^{x'}_{-\step}(y')\right)\right|}\\
 & =\frac{1}{2}\int_{y\in V_{+}}\frac{e^{-H(x',y')}}{\left|\det\left(DF^{x'}_{-\step}(y')\right)\right|}
\end{align*}
where we used that $e^{-H(x,y)}=e^{-H(x',y')}$ (Lemma \ref{lem:ham_eng_pre})
at the end. 

For the second term in (\ref{eq:ham_transit}), by the same calculation,
we have that
\[
\frac{1}{2}\int_{y\in V_{-}}\frac{e^{-H(x,y)}}{\left|\det\left(DF^{x}_{-\step}(y)\right)\right|}=\frac{1}{2}\int_{y\in V_{-}}\frac{e^{-H(x',y')}}{\left|\det\left(DF^{x'}_{\step}(y')\right)\right|}
\]
\end{proof}

\subsection*{Convergence}

First, we consider convergence when $H(x,v)=f(x)+\frac{1}{2}\norm v^{2}$
for a strongly convex function $f$. So the marginal of the stationary
distribution along $x$ is proportional to $e^{-f}$. The idea here
is coupling (as we did for Langevin dynamics). We consider two separate
processes $x$ and $y$, with their next step directions chosen to
be identical. The key lemma is that under this coupling the squared
distance decreases up to a certain time that depends on the condition
number. For a tight analysis, showing a mixing rate of $O(\kappa)$
where the condition number $\kappa$ is the ratio of the Lipschitz
constant of the gradient of $f$ to its strong convexity parameter,
we refer the reader to \cite{chen2022optimal}. The general Riemannian
setting, which gives an affine-invariant sampling algorithm, was analyzed
in \cite{LV2022geodesic,LV18riemannian}, showing that when applied
to polytopes with the Lewis weights barrier, the resulting process
converges with a mixing rate of $\tilde{O}(mn^{2/3})$ and each iteration
has time complexity $O(mn^{\omega-1})$. 
